% Tabla retirada de secciones/06-objetivos.tex (subsección «Indicadores y
% Plazos»), copiada del registro de REVISION-OBJETIVOS.md, sección 4.
% Para recuperarla: % Tabla retirada de secciones/06-objetivos.tex (subsección «Indicadores y
% Plazos»), copiada del registro de REVISION-OBJETIVOS.md, sección 4.
% Para recuperarla: % Tabla retirada de secciones/06-objetivos.tex (subsección «Indicadores y
% Plazos»), copiada del registro de REVISION-OBJETIVOS.md, sección 4.
% Para recuperarla: % Tabla retirada de secciones/06-objetivos.tex (subsección «Indicadores y
% Plazos»), copiada del registro de REVISION-OBJETIVOS.md, sección 4.
% Para recuperarla: \input{tablas/indicadores-objetivos} al final de
% 06-objetivos.tex. Usa la columna P de trabajo.tex. OJO: menciona \figarq,
% macro que ya no existe; cambiarla por Figura~\ref{fig:arquitectura}. Las
% claves OT1-OI4 corresponden a la versión anterior de los objetivos.
\subsection{Indicadores y Plazos}
La Tabla~\ref{tab:obj-smart} indica el indicador o la meta con que se comprueba
cada objetivo y la fase del Plan de Acción Preliminar (Apéndice~D) en la que
debe cumplirse. Las fechas de calendario las fijará el equipo con la profesora.

\begin{table}[tbp]
	\caption{Indicador, Evidencia y Plazo de los Objetivos}
	\label{tab:obj-smart}
	\small\linespread{1}\selectfont
	\begin{tabular}{@{}lP{7.6cm}P{3.9cm}P{2.2cm}@{}}
		\toprule
		Clave & Indicador o meta & Evidencia & Plazo (fase) \\
		\midrule
		OT1 & Diagrama de componentes y modelo entidad-relación que cubren todos los elementos de \figarq. & Diagramas & 2 \\[0.6em]
		OT2 & 100\,\% de los casos de prueba con token vencido, repetido o alterado son rechazados. & Bitácora de pruebas & 3 \\[0.6em]
		OT3 & Un ingreso validado y registrado de extremo a extremo, y tres reportes disponibles en el panel. & Demostración y registros & 2 y 3 \\[0.6em]
		OT4 & Tiempo de una lectura completa en banco (captura, decodificación y consulta) menor que el tiempo medio de validación de la línea base (Apéndice~C); 95\,\% o más de lecturas exitosas (meta provisional). & Hoja de registro de pruebas técnicas (por elaborar) & 3 (antes del piloto) \\[0.6em]
		OI1 & Tasa de llegada, tiempo medio de validación, utilización, fila máxima e incidencias calculados en los 40 intervalos de cinco días, e identificación de los intervalos con utilización igual o mayor que 1. & Hojas de observación & 1 \\[0.6em]
		OI2 & Proporción de estudiantes con celular propio, conexión y batería, y proporción que declara disposición a usar el QR, con la muestra mínima del Apéndice~A (196 a 384 estudiantes según el margen de error). & Encuestas & 1 \\[0.6em]
		OI3 & Lista de requisitos legales, cada uno con la ley de la que proviene, incorporados al diseño, con el aviso de privacidad y el periodo de conservación definidos antes del piloto. & Documento de requisitos & 1 y 2 \\[0.6em]
		OI4 & Diferencia entre línea base y piloto en tiempo medio de validación, fila máxima e incidencias, con su valor de contraste en los 40 pares de intervalos. & Hojas de observación & 4 y 5 \\
		\bottomrule
	\end{tabular}
	\tablenote{Elaboración propia. Fases del Apéndice~D: 1 diagnóstico, 2 diseño
		de hardware y firmware, 3 desarrollo del servidor, 4 piloto controlado,
		5 evaluación. Los 40 intervalos son 8 de cinco minutos por día durante
		cinco días. La meta de lecturas exitosas es una propuesta del equipo que se
		ajustará con las pruebas, y el 100\,\% de OT2 es un criterio de aceptación
		de las pruebas, no un resultado de investigación. Las incidencias de
		identidad se expresan como tasa por ingresos: (sin credencial + rechazos) /
		ingresos; un sistema que detecta mejor puede registrar más incidencias, por
		lo que se interpretan como capacidad de detección. Los casos «sin revisión»
		se reportan aparte.}
\end{table}
 al final de
% 06-objetivos.tex. Usa la columna P de trabajo.tex. OJO: menciona \figarq,
% macro que ya no existe; cambiarla por Figura~\ref{fig:arquitectura}. Las
% claves OT1-OI4 corresponden a la versión anterior de los objetivos.
\subsection{Indicadores y Plazos}
La Tabla~\ref{tab:obj-smart} indica el indicador o la meta con que se comprueba
cada objetivo y la fase del Plan de Acción Preliminar (Apéndice~D) en la que
debe cumplirse. Las fechas de calendario las fijará el equipo con la profesora.

\begin{table}[tbp]
	\caption{Indicador, Evidencia y Plazo de los Objetivos}
	\label{tab:obj-smart}
	\small\linespread{1}\selectfont
	\begin{tabular}{@{}lP{7.6cm}P{3.9cm}P{2.2cm}@{}}
		\toprule
		Clave & Indicador o meta & Evidencia & Plazo (fase) \\
		\midrule
		OT1 & Diagrama de componentes y modelo entidad-relación que cubren todos los elementos de \figarq. & Diagramas & 2 \\[0.6em]
		OT2 & 100\,\% de los casos de prueba con token vencido, repetido o alterado son rechazados. & Bitácora de pruebas & 3 \\[0.6em]
		OT3 & Un ingreso validado y registrado de extremo a extremo, y tres reportes disponibles en el panel. & Demostración y registros & 2 y 3 \\[0.6em]
		OT4 & Tiempo de una lectura completa en banco (captura, decodificación y consulta) menor que el tiempo medio de validación de la línea base (Apéndice~C); 95\,\% o más de lecturas exitosas (meta provisional). & Hoja de registro de pruebas técnicas (por elaborar) & 3 (antes del piloto) \\[0.6em]
		OI1 & Tasa de llegada, tiempo medio de validación, utilización, fila máxima e incidencias calculados en los 40 intervalos de cinco días, e identificación de los intervalos con utilización igual o mayor que 1. & Hojas de observación & 1 \\[0.6em]
		OI2 & Proporción de estudiantes con celular propio, conexión y batería, y proporción que declara disposición a usar el QR, con la muestra mínima del Apéndice~A (196 a 384 estudiantes según el margen de error). & Encuestas & 1 \\[0.6em]
		OI3 & Lista de requisitos legales, cada uno con la ley de la que proviene, incorporados al diseño, con el aviso de privacidad y el periodo de conservación definidos antes del piloto. & Documento de requisitos & 1 y 2 \\[0.6em]
		OI4 & Diferencia entre línea base y piloto en tiempo medio de validación, fila máxima e incidencias, con su valor de contraste en los 40 pares de intervalos. & Hojas de observación & 4 y 5 \\
		\bottomrule
	\end{tabular}
	\tablenote{Elaboración propia. Fases del Apéndice~D: 1 diagnóstico, 2 diseño
		de hardware y firmware, 3 desarrollo del servidor, 4 piloto controlado,
		5 evaluación. Los 40 intervalos son 8 de cinco minutos por día durante
		cinco días. La meta de lecturas exitosas es una propuesta del equipo que se
		ajustará con las pruebas, y el 100\,\% de OT2 es un criterio de aceptación
		de las pruebas, no un resultado de investigación. Las incidencias de
		identidad se expresan como tasa por ingresos: (sin credencial + rechazos) /
		ingresos; un sistema que detecta mejor puede registrar más incidencias, por
		lo que se interpretan como capacidad de detección. Los casos «sin revisión»
		se reportan aparte.}
\end{table}
 al final de
% 06-objetivos.tex. Usa la columna P de trabajo.tex. OJO: menciona \figarq,
% macro que ya no existe; cambiarla por Figura~\ref{fig:arquitectura}. Las
% claves OT1-OI4 corresponden a la versión anterior de los objetivos.
\subsection{Indicadores y Plazos}
La Tabla~\ref{tab:obj-smart} indica el indicador o la meta con que se comprueba
cada objetivo y la fase del Plan de Acción Preliminar (Apéndice~D) en la que
debe cumplirse. Las fechas de calendario las fijará el equipo con la profesora.

\begin{table}[tbp]
	\caption{Indicador, Evidencia y Plazo de los Objetivos}
	\label{tab:obj-smart}
	\small\linespread{1}\selectfont
	\begin{tabular}{@{}lP{7.6cm}P{3.9cm}P{2.2cm}@{}}
		\toprule
		Clave & Indicador o meta & Evidencia & Plazo (fase) \\
		\midrule
		OT1 & Diagrama de componentes y modelo entidad-relación que cubren todos los elementos de \figarq. & Diagramas & 2 \\[0.6em]
		OT2 & 100\,\% de los casos de prueba con token vencido, repetido o alterado son rechazados. & Bitácora de pruebas & 3 \\[0.6em]
		OT3 & Un ingreso validado y registrado de extremo a extremo, y tres reportes disponibles en el panel. & Demostración y registros & 2 y 3 \\[0.6em]
		OT4 & Tiempo de una lectura completa en banco (captura, decodificación y consulta) menor que el tiempo medio de validación de la línea base (Apéndice~C); 95\,\% o más de lecturas exitosas (meta provisional). & Hoja de registro de pruebas técnicas (por elaborar) & 3 (antes del piloto) \\[0.6em]
		OI1 & Tasa de llegada, tiempo medio de validación, utilización, fila máxima e incidencias calculados en los 40 intervalos de cinco días, e identificación de los intervalos con utilización igual o mayor que 1. & Hojas de observación & 1 \\[0.6em]
		OI2 & Proporción de estudiantes con celular propio, conexión y batería, y proporción que declara disposición a usar el QR, con la muestra mínima del Apéndice~A (196 a 384 estudiantes según el margen de error). & Encuestas & 1 \\[0.6em]
		OI3 & Lista de requisitos legales, cada uno con la ley de la que proviene, incorporados al diseño, con el aviso de privacidad y el periodo de conservación definidos antes del piloto. & Documento de requisitos & 1 y 2 \\[0.6em]
		OI4 & Diferencia entre línea base y piloto en tiempo medio de validación, fila máxima e incidencias, con su valor de contraste en los 40 pares de intervalos. & Hojas de observación & 4 y 5 \\
		\bottomrule
	\end{tabular}
	\tablenote{Elaboración propia. Fases del Apéndice~D: 1 diagnóstico, 2 diseño
		de hardware y firmware, 3 desarrollo del servidor, 4 piloto controlado,
		5 evaluación. Los 40 intervalos son 8 de cinco minutos por día durante
		cinco días. La meta de lecturas exitosas es una propuesta del equipo que se
		ajustará con las pruebas, y el 100\,\% de OT2 es un criterio de aceptación
		de las pruebas, no un resultado de investigación. Las incidencias de
		identidad se expresan como tasa por ingresos: (sin credencial + rechazos) /
		ingresos; un sistema que detecta mejor puede registrar más incidencias, por
		lo que se interpretan como capacidad de detección. Los casos «sin revisión»
		se reportan aparte.}
\end{table}
 al final de
% 06-objetivos.tex. Usa la columna P de trabajo.tex. OJO: menciona \figarq,
% macro que ya no existe; cambiarla por Figura~\ref{fig:arquitectura}. Las
% claves OT1-OI4 corresponden a la versión anterior de los objetivos.
\subsection{Indicadores y Plazos}
La Tabla~\ref{tab:obj-smart} indica el indicador o la meta con que se comprueba
cada objetivo y la fase del Plan de Acción Preliminar (Apéndice~D) en la que
debe cumplirse. Las fechas de calendario las fijará el equipo con la profesora.

\begin{table}[tbp]
	\caption{Indicador, Evidencia y Plazo de los Objetivos}
	\label{tab:obj-smart}
	\small\linespread{1}\selectfont
	\begin{tabular}{@{}lP{7.6cm}P{3.9cm}P{2.2cm}@{}}
		\toprule
		Clave & Indicador o meta & Evidencia & Plazo (fase) \\
		\midrule
		OT1 & Diagrama de componentes y modelo entidad-relación que cubren todos los elementos de \figarq. & Diagramas & 2 \\[0.6em]
		OT2 & 100\,\% de los casos de prueba con token vencido, repetido o alterado son rechazados. & Bitácora de pruebas & 3 \\[0.6em]
		OT3 & Un ingreso validado y registrado de extremo a extremo, y tres reportes disponibles en el panel. & Demostración y registros & 2 y 3 \\[0.6em]
		OT4 & Tiempo de una lectura completa en banco (captura, decodificación y consulta) menor que el tiempo medio de validación de la línea base (Apéndice~C); 95\,\% o más de lecturas exitosas (meta provisional). & Hoja de registro de pruebas técnicas (por elaborar) & 3 (antes del piloto) \\[0.6em]
		OI1 & Tasa de llegada, tiempo medio de validación, utilización, fila máxima e incidencias calculados en los 40 intervalos de cinco días, e identificación de los intervalos con utilización igual o mayor que 1. & Hojas de observación & 1 \\[0.6em]
		OI2 & Proporción de estudiantes con celular propio, conexión y batería, y proporción que declara disposición a usar el QR, con la muestra mínima del Apéndice~A (196 a 384 estudiantes según el margen de error). & Encuestas & 1 \\[0.6em]
		OI3 & Lista de requisitos legales, cada uno con la ley de la que proviene, incorporados al diseño, con el aviso de privacidad y el periodo de conservación definidos antes del piloto. & Documento de requisitos & 1 y 2 \\[0.6em]
		OI4 & Diferencia entre línea base y piloto en tiempo medio de validación, fila máxima e incidencias, con su valor de contraste en los 40 pares de intervalos. & Hojas de observación & 4 y 5 \\
		\bottomrule
	\end{tabular}
	\tablenote{Elaboración propia. Fases del Apéndice~D: 1 diagnóstico, 2 diseño
		de hardware y firmware, 3 desarrollo del servidor, 4 piloto controlado,
		5 evaluación. Los 40 intervalos son 8 de cinco minutos por día durante
		cinco días. La meta de lecturas exitosas es una propuesta del equipo que se
		ajustará con las pruebas, y el 100\,\% de OT2 es un criterio de aceptación
		de las pruebas, no un resultado de investigación. Las incidencias de
		identidad se expresan como tasa por ingresos: (sin credencial + rechazos) /
		ingresos; un sistema que detecta mejor puede registrar más incidencias, por
		lo que se interpretan como capacidad de detección. Los casos «sin revisión»
		se reportan aparte.}
\end{table}
